\begin{abstract}
Four-bit Microscaling (MX) formats reach acceptable accuracy for many
inference workloads, but some models and deployment targets demand more
precision than a single four-bit block provides, without conceding the memory
advantage that motivates four-bit weights in the first place. Several augmented
MX variants address this at the format level; their hardware cost, however, is
seldom measured against the alternatives they displace on common silicon.

This work introduces \fpfourres, an MX variant that represents activations as
the sum of two \fpfour{} blocks under independent block scales while keeping
weights in plain \fpfour, and evaluates whether it merits a hardware
implementation. Because only activations carry the second block, weight
footprint remains that of \fpfour{} while representational precision approaches
that of \fpeight. The value of this trade depends on delivered throughput,
which a format-level analysis cannot establish.

We implement \fpfourres{} alongside \fpfour, \fpeight, and \mtwo{} as
format-specialised independent dot-product engines within a single \xif-compliant
coprocessor (CV32E40X), characterising all four on one host core, one synthesis flow, and
one technology node. Each engine accumulates in a scale-free frame whose width
depends on the element format alone.

Synthesised at \SI{45}{\nano\meter}, the specialised \fpfour{} engine sustains
\SI{5.1}{\times} the throughput per unit area of \fpeight. The \fpfourres{}
datapath requires three source operands and two block scales, exceeding a
three-read-port register file, and is expressed as an instruction pair whose
present throughput is \DATAGAP{X}$\times$ that of \fpeight; its advantage
therefore lies in halved weight memory at comparable accuracy rather than in
compute.
\end{abstract}
